\begin{textsample}{POS Dim 1 – vem\_ed\_10 – Score 14.00 – vem\_ed\_10/t040.txt}  \label{ex:f1_pos_117}
Jornada de PCIs: abertura Em 26 de novembro de 2021, das 14 às 17h, a equipe dos Projetos Colaborativos Internacionais (PCIs / Cesu / CPS) realizou a Jornada de Projetos Colaborativos Internacionais (PCIs / Cesu): Reflexões sobre as Práticas.
O evento online \textbf{contou} com o \textbf{apoio} da Fatec Guaratinguetá e a participação das Fatecs Americana, Baixada Santista, Bauru, Garça, Guaratinguetá, Itaquaquecetuba, Itu, Lins, Piracicaba, São José do Rio Preto e São Caetano do Sul.
A gravação do webinário encontra-se em https://www.youtube.com/watch?v=acsdEqNgLPA.
As discussões se organizaram em uma abertura e quatro painéis temáticos, com apresentação de Regiane Camargo, da equipe dos PCIs, e professora da Fatec Guaratinguetá.
Da abertura, participaram Rafael Ferreira Alves (coordenador da Unidade do Ensino Superior de Graduação / Cesu), André Luiz Braun Galvão (diretor do Grupo Acadêmico Pedagógico / Cesu), André Ricardo Soares Amarante (diretor da Fatec Guaratinguetá), William Menezes (gestor pedagógico regional / Cesu), Mariane Teixeira (coordenadora de Línguas - Cesu), Osvaldo Succi Junior (coordenador dos PCIs / Cesu), Daives Bergamasco e Neusa Haruka Sezaki Gritti (equipe dos PCIs).
Rafael Ferreira Alves deu as boas-vindas e comentou o crescimento exponencial dos PCIs nos últimos anos.
André Braun ressaltou a \textbf{importância} do trabalho \textbf{colaborativo} e do \textbf{aprendizado} por projetos, dentro do \textbf{contexto} das metodologias \textbf{ativas}.
Destacou ainda o \textbf{desenvolvimento} de \textbf{competências} \textbf{socioemocionais} de professores e alunos.
André Amarante agradeceu os esforços da Cesu e frisou a \textbf{importância} dos PCIs na formação dos egressos das Fatecs.
Mariane Teixeira ressaltou \textbf{que} a realização dos PCIs contribui para a internacionalização do currículo dos cursos superiores de tecnologia e agradeceu a participação dos professores de línguas para o sucesso dos projetos.
Osvaldo Succi comentou sobre as iniciativas de internacionalização desenvolvida nas Fatecs e \textbf{mencionou} o \textbf{apoio} da equipe dos PCIs / Cesu à ArInter e à Cetec na realização do Procin, para internacionalização do ensino médio e técnico.
Neusa Haruka apresentou a equipe dos PCIs.
Daives Bergamasco \textbf{mencionou} a \textbf{importância} das plataformas para atividades \textbf{síncronas} (Zoom, Skype) e assíncronas (Slack, Lark).

% matched lemmas: apoio, aprendizado, ativo, colaborativo, competência, contar, contexto, desenvolvimento, importância, mencionar, que, socioemocional, síncrono
\end{textsample}
